% Преамбула: шрифты (латиница, кириллица, арабский с огласовками), поля, библиография.
\usepackage{etoolbox}
\usepackage{fontspec}
\usepackage{polyglossia}
\setdefaultlanguage{\mainlang}
\ifdefstring{\mainlang}{english}{\setotherlanguages{russian,arabic}}{\setotherlanguages{english,arabic}}

% Times New Roman есть в macOS; на другой машине — Liberation Serif или Tempora.
\setmainfont{Times New Roman}
\newfontfamily\cyrillicfont{Times New Roman}
% Арабский — тот же шрифт, что рисует карточки бота (ADR-0024): огласовки и шадда не теряются.
\newfontfamily\arabicfont[Script=Arabic,Path=../src/muallim/presentation/telegram/assets/fonts/,
  Extension=.ttf,UprightFont=*-Regular,BoldFont=*-Bold,Scale=1.25]{ScheherazadeNew}
% \ar{…} — арабская вставка в строке; источник строки — в тексте рядом ([SRC-NNN]).
\newcommand{\ar}[1]{\textarabic{#1}}

\usepackage[left=3cm,right=1.5cm,top=2cm,bottom=2cm]{geometry} % [утвердить] по требованиям вуза
\usepackage{setspace}
\onehalfspacing
\usepackage{indentfirst}

\usepackage{graphicx}
\graphicspath{{../docs/architecture/diagrams/}{../docs/design/screens/}{figures/}}
\usepackage{booktabs,tabularx,longtable}
\usepackage{caption}
\usepackage{csquotes}
\usepackage[hidelinks]{hyperref}

% Библиография — тот же sources.bib, что у всего проекта (ключи SRC-NNN).
\ifdefstring{\mainlang}{english}
  {\usepackage[backend=biber,style=ieee,sorting=none]{biblatex}}
  {\usepackage[backend=biber,style=gost-numeric,sorting=none]{biblatex}}
\addbibresource{../docs/research/sources.bib}

% ID рисунков и таблиц из docs/thesis/figures.md: \label{fig:R-01}, \label{tab:T-01}.
